\section{Additional Proofs}
\label{app:proofs}

\subsection{Polynomial and Analytic Approximation}

\begin{proof}[Proof of Theorem~\ref{thm:poly-generation}]
Since \(g\) has degree at most \(P\), write
\begin{equation}
  g(s)=\sum_{p=0}^{P} c_p s^p.
\end{equation}
Take the NTD-PL latent tensor to be the true latent tensor,
\(\cS=\cS^\star\), and set \(\beta_p=c_p\). Then, entrywise,
\begin{equation}
  \widehat{\cY}_{i_1,\ldots,i_N}
  =
  \sum_{p=0}^{P} \beta_p
  \left(\cS^\star_{i_1,\ldots,i_N}\right)^p
  =
  g\left(\cS^\star_{i_1,\ldots,i_N}\right)
  =
  \cY_{i_1,\ldots,i_N}.
\end{equation}
Thus \(\widehat{\cY}=\cY\).
\end{proof}

\begin{lemma}[Scalar-to-tensor approximation]
\label{lem:scalar-to-tensor}
Let \(\norm{\cS^\star}_\infty\le B\). For any scalar function \(g\) and any
degree-\(P\) polynomial \(q_P\),
\begin{equation}
  \norm{g(\cS^\star)-q_P(\cS^\star)}_F
  \le
  \sqrt{\prod_{n=1}^{N}I_n}\,
  \sup_{|s|\le B}|g(s)-q_P(s)|.
\end{equation}
\end{lemma}

\begin{proof}
Every entry of \(\cS^\star\) lies in \([-B,B]\), so every entrywise
approximation error is bounded by
\(\sup_{|s|\le B}|g(s)-q_P(s)|\). Squaring, summing over all entries, and
taking the square root gives the claim.
\end{proof}

\begin{proof}[Proof of Theorem~\ref{thm:analytic-approx}]
Let \(g(s)=\sum_{k=0}^{\infty} a_k s^k\) be the Taylor expansion of \(g\)
around zero, and choose \(q_P(s)=\sum_{k=0}^{P}a_k s^k\). By Cauchy's
estimate,
\begin{equation}
  |a_k| \le \frac{M_\rho}{\rho^k}.
\end{equation}
For \(|s|\le B<\rho\),
\begin{align}
  |g(s)-q_P(s)|
  &\le
  \sum_{k=P+1}^{\infty}|a_k|\,|s|^k \\
  &\le
  M_\rho\sum_{k=P+1}^{\infty}(B/\rho)^k
  =
  M_\rho\frac{(B/\rho)^{P+1}}{1-B/\rho}.
\end{align}
Applying Lemma~\ref{lem:scalar-to-tensor} and taking
\(\cS=\cS^\star\), \(\beta_k=a_k\) for \(0\le k\le P\), yields the stated
unnormalized bound. Dividing by
\(\sqrt{M}\), \(M=\prod_{n=1}^{N}I_n\), gives the normalized form used in
Theorem~\ref{thm:analytic-approx}.
\end{proof}

\subsection{Rank Expansion}

\begin{proof}[Proof of Proposition~\ref{prop:rank-expansion}]
Write the Tucker tensor entrywise as
\begin{equation}
  \cS_{i_1,\ldots,i_N}
  =
  \sum_{r_1=1}^{R_1}\cdots\sum_{r_N=1}^{R_N}
  \cX_{r_1,\ldots,r_N}
  \prod_{n=1}^{N} A^{(n)}_{i_n,r_n}.
\end{equation}
Taking an element-wise \(p\)-th power gives
\begin{align}
  \left(\cS_{i_1,\ldots,i_N}\right)^p
  &=
  \prod_{q=1}^{p}
  \left(
  \sum_{r_1^{(q)}=1}^{R_1}\cdots
  \sum_{r_N^{(q)}=1}^{R_N}
  \cX_{r_1^{(q)},\ldots,r_N^{(q)}}
  \prod_{n=1}^{N} A^{(n)}_{i_n,r_n^{(q)}}
  \right) \\
  &=
  \sum_{\gamma_1\in[R_1]^p}\cdots
  \sum_{\gamma_N\in[R_N]^p}
  \left(
  \prod_{q=1}^{p}
  \cX_{r_1^{(q)},\ldots,r_N^{(q)}}
  \right)
  \prod_{n=1}^{N}
  \left(
  \prod_{q=1}^{p} A^{(n)}_{i_n,r_n^{(q)}}
  \right),
\end{align}
where \(\gamma_n=(r_n^{(1)},\ldots,r_n^{(p)})\). For each mode \(n\), let
\(\alpha_n\in\mathbb N^{R_n}\) record the counts of the entries in
\(\gamma_n\), so \(|\alpha_n|=p\). Then
\begin{equation}
  \prod_{q=1}^{p} A^{(n)}_{i_n,r_n^{(q)}}
  =
  \prod_{r=1}^{R_n}\left(A^{(n)}_{i_n,r}\right)^{\alpha_{n,r}}
  =
  A^{(n)\langle p\rangle}_{i_n,\alpha_n}.
\end{equation}
Group all ordered tuples \((\gamma_1,\ldots,\gamma_N)\) that produce the same
multi-indices \((\alpha_1,\ldots,\alpha_N)\), and define the aggregated core
\begin{equation}
  \cX^{\langle p\rangle}_{\alpha_1,\ldots,\alpha_N}
  =
  \sum_{\substack{\gamma_1,\ldots,\gamma_N:\\
  \operatorname{count}(\gamma_n)=\alpha_n\ \forall n}}
  \prod_{q=1}^{p}
  \cX_{r_1^{(q)},\ldots,r_N^{(q)}}.
\end{equation}
The preceding display is exactly the Tucker reconstruction
\begin{equation}
  \cS^{\odot p}
  =
  \tucker{\cX^{\langle p\rangle}}
  {\bA^{(1)\langle p\rangle},\ldots,\bA^{(N)\langle p\rangle}}.
\end{equation}
The \(n\)-th Veronese factor has \(\binom{R_n+p-1}{p}\) columns, giving the
claimed Tucker rank bound.
\end{proof}

\subsection{Separation}

\begin{lemma}[Full-rank Veronese evaluations]
\label{lem:veronese-full-rank}
Let \(D=\binom{R+p-1}{p}\). There exists a matrix \(\bA\in\R^{D\times R}\)
whose Veronese feature matrix \(\bA^{\langle p\rangle}\in\R^{D\times D}\)
has rank \(D\). Consequently, for \(I\ge D\), a generic matrix
\(\bA\in\R^{I\times R}\) has \(\operatorname{rank}(\bA^{\langle p\rangle})=D\).
\end{lemma}

\begin{proof}
The homogeneous degree-\(p\) monomials in \(R\) variables are linearly
independent polynomials. Hence there exist evaluation points
\(x_1,\ldots,x_D\in\R^R\) such that the matrix
\((x_i^\alpha)_{i,\alpha}\), \(|\alpha|=p\), is nonsingular; otherwise every
\(D\times D\) evaluation determinant would vanish, which would imply a
nontrivial linear dependence among the monomials. Taking these points as the
rows of \(\bA\) gives the first claim.

For \(I\ge D\), choose any \(D\) rows. The determinant of the corresponding
\(D\times D\) Veronese submatrix is a polynomial in the entries of \(\bA\)
that is not identically zero by the first claim. Its nonvanishing is therefore
a generic property, so \(\bA^{\langle p\rangle}\) has full column rank for
generic \(\bA\).
\end{proof}

\begin{proof}[Proof of Theorem~\ref{thm:tucker-separation}]
Let \(D=\binom{R+p-1}{p}\). Choose generic factor matrices
\(\bA^{(n)}=[a_1^{(n)},\ldots,a_R^{(n)}]\in\R^{I_n\times R}\), and define
the CP-diagonal Tucker backbone
\begin{equation}
  \cS
  =
  \sum_{r=1}^{R}
  a_r^{(1)}\otimes\cdots\otimes a_r^{(N)}.
\end{equation}
Its mode-\(n\) unfolding is
\begin{equation}
  \cS_{(n)}
  =
  \bA^{(n)}
  \left(
  \bA^{(N)}\odot\cdots\odot\bA^{(n+1)}
  \odot\bA^{(n-1)}\odot\cdots\odot\bA^{(1)}
  \right)^\top .
\end{equation}
For generic \(\bA^{(n)}\), \(\bA^{(n)}\) has rank \(R\), and the Khatri-Rao
product on the right has rank \(R\). Thus
\(\operatorname{rank}(\cS_{(n)})=R\) for every \(n\), so \(\cS\) has exact
Tucker rank \((R,\ldots,R)\).

Set \(\beta_p=1\) and \(\beta_q=0\) for \(q\ne p\). Then
\(\cY=\cS^{\odot p}\) is a valid degree-\(p\) NTD-PL tensor. Entrywise,
\begin{align}
  \cY_{i_1,\ldots,i_N}
  &=
  \left(
  \sum_{r=1}^{R}\prod_{n=1}^{N} A^{(n)}_{i_n,r}
  \right)^p \\
  &=
  \sum_{|\alpha|=p}
  c_\alpha
  \prod_{n=1}^{N}\prod_{r=1}^{R}
  \left(A^{(n)}_{i_n,r}\right)^{\alpha_r},
\end{align}
where \(c_\alpha=p!/(\alpha_1!\cdots\alpha_R!)>0\). Let
\(\bB^{(n)}=\bA^{(n)\langle p\rangle}\in\R^{I_n\times D}\). By
Lemma~\ref{lem:veronese-full-rank}, each \(\bB^{(n)}\) has full column rank
\(D\) for generic \(\bA^{(n)}\). The mode-\(n\) unfolding is
\begin{equation}
  \cY_{(n)}
  =
  \bB^{(n)}\operatorname{diag}(c)
  \left(
  \bB^{(N)}\odot\cdots\odot\bB^{(n+1)}
  \odot\bB^{(n-1)}\odot\cdots\odot\bB^{(1)}
  \right)^\top .
\end{equation}
The diagonal matrix \(\operatorname{diag}(c)\) is nonsingular. Since a
Khatri-Rao product of full-column-rank matrices is full column rank, the
right factor has rank \(D\). Therefore \(\operatorname{rank}(\cY_{(n)})=D\).
This holds for every mode \(n\), and the exact Tucker rank of \(\cY\) is
\((D,\ldots,D)\).
\end{proof}

\subsection{Masked Generalization}
\label{app:masked-generalization}

\begin{theorem}[Masked-entry uniform convergence]
\label{thm:masked-generalization}
Let \(\Theta\subset\R^d\) be a bounded NTD-PL parameter set with
\[
  d=
  \prod_{n=1}^{N}R_n
  +
  \sum_{n=1}^{N}I_nR_n
  +
  (P+1).
\]
For \(\theta\in\Theta\), define the full-entry risk
\[
  \mathcal L(\theta)
  =
  \frac{1}{M}\sum_{\mathbf i}\ell_{\mathbf i}(\theta),
  \qquad
  M=\prod_{n=1}^{N}I_n,
\]
and the masked empirical risk
\[
  \widehat{\mathcal L}_{m}(\theta)
  =
  \frac{1}{m}\sum_{t=1}^{m}
  \ell_{\mathbf i_t}(\theta),
\]
where \(\mathbf i_t\) are sampled uniformly from tensor entries. Assume
\(0\le \ell_{\mathbf i}(\theta)\le B_\ell\) and
\(\ell_{\mathbf i}(\theta)\) is \(G\)-Lipschitz in \(\theta\), uniformly over
\(\mathbf i\). If \(\Theta\) is contained in a Euclidean ball of radius
\(B_\Theta\), then with probability at least \(1-\delta\),
\[
  \sup_{\theta\in\Theta}
  \left|
  \mathcal L(\theta)-\widehat{\mathcal L}_{m}(\theta)
  \right|
  \le
  2G\varepsilon
  +
  B_\ell
  \sqrt{
  \frac{
  d\log(3B_\Theta/\varepsilon)+\log(2/\delta)}
  {2m}}
\]
for every \(0<\varepsilon\le B_\Theta\).
\end{theorem}

\begin{proof}[Proof of Theorem~\ref{thm:masked-generalization}]
Fix \(0<\varepsilon\le B_\Theta\), and let
\(\mathcal N_\varepsilon\) be an \(\varepsilon\)-net of \(\Theta\) in
Euclidean norm. Since \(\Theta\) lies in a \(d\)-dimensional ball of radius
\(B_\Theta\), it admits a net with
\begin{equation}
  |\mathcal N_\varepsilon|
  \le
  \left(\frac{3B_\Theta}{\varepsilon}\right)^d.
\end{equation}
For any fixed \(\theta_0\in\mathcal N_\varepsilon\), the random variables
\(\ell_{\mathbf i_t}(\theta_0)\) are independent and bounded in
\([0,B_\ell]\). Hoeffding's inequality gives
\begin{equation}
  \Pr\left(
  \left|
  \mathcal L(\theta_0)-\widehat{\mathcal L}_m(\theta_0)
  \right|
  \ge u
  \right)
  \le
  2\exp\left(-\frac{2mu^2}{B_\ell^2}\right).
\end{equation}
Taking a union bound over \(\mathcal N_\varepsilon\), with probability at
least \(1-\delta\),
\begin{equation}
  \max_{\theta_0\in\mathcal N_\varepsilon}
  \left|
  \mathcal L(\theta_0)-\widehat{\mathcal L}_m(\theta_0)
  \right|
  \le
  B_\ell
  \sqrt{
  \frac{
  \log(2|\mathcal N_\varepsilon|/\delta)}
  {2m}}.
\end{equation}
Substituting the net-size bound yields the stated stochastic term.

Now take any \(\theta\in\Theta\), and choose
\(\theta_0\in\mathcal N_\varepsilon\) with
\(\norm{\theta-\theta_0}_2\le\varepsilon\). By the uniform Lipschitz
assumption,
\begin{equation}
  |\mathcal L(\theta)-\mathcal L(\theta_0)|
  \le G\varepsilon,
  \qquad
  |\widehat{\mathcal L}_m(\theta)-\widehat{\mathcal L}_m(\theta_0)|
  \le G\varepsilon.
\end{equation}
Therefore
\begin{align}
  \left|
  \mathcal L(\theta)-\widehat{\mathcal L}_m(\theta)
  \right|
  &\le
  |\mathcal L(\theta)-\mathcal L(\theta_0)|
  +
  \left|
  \mathcal L(\theta_0)-\widehat{\mathcal L}_m(\theta_0)
  \right| \\
  &\quad+
  |\widehat{\mathcal L}_m(\theta_0)-\widehat{\mathcal L}_m(\theta)| \\
  &\le
  2G\varepsilon
  +
  B_\ell
  \sqrt{
  \frac{
  d\log(3B_\Theta/\varepsilon)+\log(2/\delta)}
  {2m}}.
\end{align}
Taking the supremum over \(\theta\in\Theta\) completes the proof.
\end{proof}

\begin{remark}
The Lipschitz and bounded-loss assumptions are standard for compact
parametric classes. For NTD-PL, they follow when the observed values, core,
factor matrices, and polynomial coefficients are bounded and the polynomial
degree \(P\) is finite. The resulting Lipschitz constant may depend
polynomially on these bounds and on \(P\), but the covering dimension remains
the number of tied NTD-PL parameters rather than the number of parameters in
an untied high-order Tucker expansion.
\end{remark}
